\footnotesize%
\pstart%
\noindent%
In der dichten Aufzeichnung N.\;\ref{57037-1} \textit{Memorabile inventum barometri et thermometri portatilis exactissimi}, welche Leibniz selbst auf August 1677 datiert hat, wird in Anlehnung an Überlegungen von Christiaan Huygens\indexN{\textls{Huygens} (Hugenius, Ugenius, Hugens, Huguens), Christiaan 1629\textendash1695} und Louis Hubin\indexN{\textls{Hubin} (Hubinus), Louis (Ludion de) fl. 1673\textendash1683} ein Instrument umrissen, das die physikalischen Eigenschaften eines Barometers und eines Thermometers vereinigen und zugleich den Vorteil der Tragbarkeit aufweisen soll.
Dieser Umriss wurde einige Wochen später in zwei aufeinander folgenden Aufzeichnungen \textit{De barometro portabili} (N.\;\ref{57038} und N.\;\ref{57039}) wiederaufgenommen und weiter entwickelt;
beide \textit{Schedae} wurden von Leibniz eigenhändig auf den 12.\;(22.) November 1677 datiert.
Insbesondere die Aufzeichnung N.\;\ref{57038} setzt den Kerngedanken von N.\;\ref{57037-1} fort und entfaltet ihn, während die Aufzeichnung N.\;\ref{57039} nach dem gestrichenen, an N.\;\ref{57038} unmittelbar anknüpfenden Anfang sich Überlegungen über das Verhältnis zweiter Fluiden (Quecksilber und Luft) in einem Barometer widmet.
\pend%
\pstart%
Der textgenetische und inhaltliche Zusammenhang der drei Aufzeichnungen N.\;\ref{57037-1}, \ref{57038} und \ref{57039} rechfertigt deren einheitlichen Edition im Textkomplex N.\;\ref{barom-portab_intro} \textit{Barometrum portabile}.
Die Gesamtdatierung von N.\;\ref{barom-portab_intro} ist durch die eigenhändigen Datierungen der einzelnen Aufzeichnungen festgelegt.
\pend%
\selectlanguage{latin}%
\frenchspacing%
\normalsize%
